\documentclass{beamer}

% Theme — deliberately different from Madrid (rounded header/footer bars):
% Boadilla uses plain horizontal-rule frame titles and a slim footline,
% no boxed title page, no shadowed blocks.
\usetheme{Boadilla}
\usecolortheme{orchid}
\usefonttheme{professionalfonts}
\setbeamertemplate{navigation symbols}{}
\setbeamertemplate{blocks}[rounded][shadow=false]
\usepackage{tikz}

% Packages
\usepackage{graphicx}
\usepackage{amsmath}
\usepackage{amssymb}

%-------------------------------------------------
% Title Page Details
%-------------------------------------------------
\title{\textbf{Mathematics for Computing}}
\subtitle{Assignment}

\author{%
Aleena Mary John 
Department of Computer Science and Engineering with Artificial Intelligence and Software Engineering \\[0.2cm]
}

\institute{
CUSAT
}

\date{\today}

%-------------------------------------------------
\begin{document}

%-------------------------------------------------
% Title Slide
%-------------------------------------------------
\begin{frame}
    \titlepage
\end{frame}

%-------------------------------------------------
% Content
%-------------------------------------------------
\begin{frame}{Content}
\begin{itemize}
    \item Set-Builder Form (10 Examples)
    \item Types of Sets with AI/DS Examples
    \item Inclusion-Exclusion Principle
    \item Hasse Diagram of a Poset
\end{itemize}
\end{frame}

%-------------------------------------------------
% Slide: Set-Builder Form (Part 1)
%-------------------------------------------------
\begin{frame}{Set-Builder Form --- Examples (1--5)}

\begin{enumerate}
    \item Even numbers:
    \[
    \{x \mid x = 2n,\; n \in \mathbb{Z}\}
    \]
    \item Natural numbers less than 10:
    \[
    \{x \mid x \in \mathbb{N},\; x < 10\}
    \]
    \item Vowels in English:
    \[
    \{x \mid x \text{ is a vowel in the English alphabet}\}
    \]
    \item Perfect squares up to 100:
    \[
    \{x \mid x = n^2,\; n \in \mathbb{Z},\; 1 \leq n \leq 10\}
    \]
    \item Odd integers:
    \[
    \{x \mid x \in \mathbb{Z},\; x = 2k+1,\; k \in \mathbb{Z}\}
    \]
\end{enumerate}

\end{frame}

%-------------------------------------------------
% Slide: Set-Builder Form (Part 2)
%-------------------------------------------------
\begin{frame}{Set-Builder Form --- Examples (6--10)}

\begin{enumerate}
    \setcounter{enumi}{5}
    \item Primes less than 20:
    \[
    \{x \mid x \text{ is prime},\; x < 20\}
    \]
    \item Multiples of 5 between 1 and 50:
    \[
    \{x \mid x = 5n,\; n \in \mathbb{Z},\; 1 \leq 5n \leq 50\}
    \]
    \item Reals between 0 and 1:
    \[
    \{x \mid x \in \mathbb{R},\; 0 < x < 1\}
    \]
    \item Days of the week:
    \[
    \{x \mid x \text{ is a day of the week}\}
    \]
    \item Solutions of $x^2 - 4 = 0$:
    \[
    \{x \mid x \in \mathbb{R},\; x^2 - 4 = 0\} = \{-2, 2\}
    \]
\end{enumerate}

\end{frame}

%-------------------------------------------------
% Slide: Types of Sets (Part 1)
%-------------------------------------------------
\begin{frame}{Types of Sets --- AI/DS Examples (I)}

\begin{itemize}
    \item \textbf{Empty (Null) Set} --- has no elements $(\emptyset)$. \\
    \textit{Predictions with model confidence $>100\%$.}
    \vspace{0.15cm}
    \item \textbf{Singleton Set} --- exactly one element. \\
    \textit{Unique optimal weight vector for a strictly convex loss.}
    \vspace{0.15cm}
    \item \textbf{Finite Set} --- countable, ends. \\
    \textit{Class labels: \{spam, not spam\}.}
    \vspace{0.15cm}
    \item \textbf{Infinite Set} --- doesn't end. \\
    \textit{All possible real-valued weights of a neuron.}
    \vspace{0.15cm}
    \item \textbf{Equal Sets} --- same elements, any order. \\
    \textit{Two feature-selection runs agreeing on selected features.}
\end{itemize}

\end{frame}

%-------------------------------------------------
% Slide: Types of Sets (Part 2)
%-------------------------------------------------
\begin{frame}{Types of Sets --- AI/DS Examples (II)}

\begin{itemize}
    \item \textbf{Equivalent Sets} --- same cardinality. \\
    \textit{Train and validation sets, both of size 1{,}000.}
    \vspace{0.15cm}
    \item \textbf{Subset} --- every element of $A$ is in $B$. \\
    \textit{Positive predictions $\subseteq$ all predictions.}
    \vspace{0.15cm}
    \item \textbf{Universal Set} --- contains everything relevant. \\
    \textit{Full vocabulary in an NLP tokenizer.}
    \vspace{0.15cm}
    \item \textbf{Power Set} --- set of all subsets. \\
    \textit{All candidate feature subsets in wrapper-method selection.}
    \vspace{0.15cm}
    \item \textbf{Disjoint Sets} --- no common elements. \\
    \textit{Training set and test set (properly split, no overlap).}
\end{itemize}

\end{frame}

%-------------------------------------------------
% Slide: Inclusion-Exclusion Principle
%-------------------------------------------------
\begin{frame}{Inclusion--Exclusion Principle}

\textbf{For two sets:}
\[
|A \cup B| = |A| + |B| - |A \cap B|
\]
Add both sets, then subtract the overlap once, since it was counted twice.

\vspace{0.4cm}

\textbf{For three sets:}
\[
|A\cup B\cup C| = |A|+|B|+|C|-|A\cap B|-|A\cap C|-|B\cap C|+|A\cap B\cap C|
\]

\vspace{0.4cm}

\textbf{General pattern:} add all single sets, subtract all pairwise
intersections, add back all triple intersections, and so on ---
alternating signs as you go deeper.

\vspace{0.3cm}

\textbf{DS use case:} counting users who clicked at least one of three
ad campaigns, given the overlaps between campaigns.

\end{frame}

%-------------------------------------------------
% Slide: Hasse Diagram
%-------------------------------------------------
\begin{frame}{Hasse Diagram of a Poset}

\textbf{Poset:} a set with a relation that is reflexive, antisymmetric,
and transitive --- not every pair of elements needs to be comparable.

\vspace{0.2cm}

\textbf{Hasse Diagram:} draws only the covering relations (direct,
non-redundant ``less than'' links), with greater elements placed higher.
No arrows or self-loops are needed.

\vspace{0.4cm}

\textbf{Example: } $(\{1,2,3,6\}, \mid)$ --- ``divides''

\begin{center}
\begin{tikzpicture}[scale=1.2]

% Nodes
\node (6) at (0,2) {6};
\node (2) at (-1,1) {2};
\node (3) at (1,1) {3};
\node (1) at (0,0) {1};

% Edges
\draw (1)--(2);
\draw (1)--(3);
\draw (2)--(6);
\draw (3)--(6);

\end{tikzpicture}
\end{center}

No line is drawn directly from 1 to 6, since that relation is already
implied by $1 \to 2 \to 6$ and $1 \to 3 \to 6$.

\end{frame}

\end{document}